\documentclass[hidelinks]{article}
\usepackage[a4paper, total={7in, 10in}]{geometry}
\usepackage[dvipsnames]{xcolor}
\usepackage{amsmath}
\usepackage{tikz}
\usepackage{tkz-euclide}
\usepackage[unicode]{hyperref}
\usepackage[all]{hypcap}
\usepackage{fancyhdr}
\usepackage{tocloft} % Add this package
\usepackage{xcolor}
\definecolor{darkgreen}{RGB}{0,100,0}
\usepackage{listings}
\lstset{
  language=R,
  basicstyle=\ttfamily\small,
  keywordstyle=\color{red},
  commentstyle=\color{darkgreen},
  stringstyle=\color{darkgreen},
  numbers=left,
  numberstyle=\tiny,
  stepnumber=1,
  numbersep=5pt,
  frame=single,
  breaklines=true,
  showstringspaces=false
}
\usepackage{amsmath}
\usepackage{amssymb}
\usepackage{pgfplots}
\pgfplotsset{compat=1.18}
\usepackage{amsmath}
\usepackage{float}
\usepackage{amsthm}
\usepackage[shortlabels]{enumitem}



% Remove dots from table of contents
\renewcommand{\cftdot}{}
\renewcommand{\cftsecdotsep}{\cftnodots}
\renewcommand{\cftsubsecdotsep}{\cftnodots}
\newcommand{\RNum}[1]{\uppercase\expandafter{\romannumeral #1\relax}}


\renewcommand{\contentsname}{Table of Contents}
\usetikzlibrary{angles,calc, decorations.pathreplacing}
\definecolor{carminered}{rgb}{1.0, 0.0, 0.22}
\definecolor{capri}{rgb}{0.0, 0.75, 1.0}
\definecolor{brightlavender}{rgb}{0.75, 0.58, 0.89}

\title{\textbf{STAT 545 HW 4}}
\author{Andres Efren Rocha Jayasinha}
\setcounter{secnumdepth}{0}
\date{Updated: \today} % suppresses the printed date

\begin{document}
\hypersetup{bookmarksnumbered=true,}
\pagecolor{white}
\color{black}
\maketitle
\begin{Large}
\tableofcontents
\end{Large}%
\pagebreak

\section{Assignment 4: Metropolis-Hastings algorithm}

\subsection{Exercise 5.7}

Exhibit a Metropolis-Hastings algorithm to sample from a binomial distribution with parameters n and p. Use a proportional distribution that is uniform on \{0,1,...,n\}.

\smallskip

\noindent \textbf{Solution}

\noindent We wish to construct a Metropolis--Hastings (MH) algorithm to sample from a 
Binomial distribution with parameters $n$ and $p$, whose probability mass 
function is

\[
\pi(k) = \mathbb{P}(X = k)
= \binom{n}{k} p^k (1-p)^{n-k},
\qquad k \in \{0,1,\dots,n\}.
\]

\noindent We are instructed to use a proposal distribution that is uniform on 
$\{0,1,\dots,n\}$. Thus, for any current state $x$,

\[
q(y \mid x) = \frac{1}{n+1}, 
\qquad y \in \{0,1,\dots,n\}.
\]

\noindent This is an independence proposal (it does not depend on $x$), and it is symmetric in the sense that

\[
q(x \mid y) = q(y \mid x) = \frac{1}{n+1}.
\]


\noindent The general Metropolis--Hastings acceptance probability is

\[
\alpha(x,y)
=
\min\left(
1,
\frac{\pi(y)\, q(x \mid y)}{\pi(x)\, q(y \mid x)}
\right).
\]

\noindent Since the proposal is symmetric, $q(x \mid y) = q(y \mid x)$ and these terms cancel. Therefore,

\[
\alpha(x,y)
=
\min\left(
1,
\frac{\pi(y)}{\pi(x)}
\right).
\]

\noindent Substituting the Binomial pmf,

\[
\frac{\pi(y)}{\pi(x)}
=
\frac{
\binom{n}{y} p^y (1-p)^{n-y}
}{
\binom{n}{x} p^x (1-p)^{n-x}
}
=
\frac{\binom{n}{y}}{\binom{n}{x}}
\left( \frac{p}{1-p} \right)^{y-x}.
\]

\noindent Hence the acceptance probability is

\[
\boxed{
\alpha(x,y)
=
\min\left(
1,
\frac{\binom{n}{y}}{\binom{n}{x}}
\left( \frac{p}{1-p} \right)^{y-x}
\right).
}
\]

\noindent Let $T$ be the total number of iterations, and choose an initial value 
$x_0 \in \{0,1,\dots,n\}$.

\noindent For $t = 1,2,\dots,T$:

\begin{enumerate}
    \item Propose $y \sim \text{Unif}\{0,1,\dots,n\}$.
    \item Compute
    \[
    \alpha
    =
    \min\left(
    1,
    \frac{\pi(y)}{\pi(x_{t-1})}
    \right).
    \]
    \item Draw $u \sim \text{Unif}(0,1)$.
    \item If $u \le \alpha$, set $x_t = y$ (accept).
    \item Otherwise, set $x_t = x_{t-1}$ (reject).
\end{enumerate}

\noindent The Markov chain $\{x_t\}$ constructed in this way has stationary 
distribution $\text{Binomial}(n,p)$. After discarding an initial burn-in 
period, the samples may be used as approximate draws from the target 
distribution.

\subsection{Exercise 5.8}

The Metropolis--Hastings algorithm is used to simulate a binomial random variable with parameters
$n=4$ and $p=\tfrac14$. The proposal distribution is a simple symmetric random walk on
$\{0,1,2,3,4\}$ with reflecting boundaries.
\begin{enumerate}
  \item[(a)] Exhibit the $T$ matrix.
  \item[(b)] Exhibit the transition matrix for the Markov chain created by the Metropolis--Hastings algorithm.
  \item[(c)] Show explicitly that this transition matrix gives a reversible Markov chain whose stationary
  distribution is the desired binomial distribution.
\end{enumerate}

\noindent \textbf{Solution}

\noindent We want to simulate the target distribution
\[
\pi(k)=\binom{4}{k}\left(\frac14\right)^k\left(\frac34\right)^{4-k},\qquad k\in\{0,1,2,3,4\},
\]
using a Metropolis--Hastings (MH) chain whose proposal is a simple symmetric random walk on
\(\{0,1,2,3,4\}\) with reflecting boundaries.

\smallskip

\noindent Part A:

\noindent A simple symmetric random walk proposes moving one step left or right with probability \(1/2\)
when in an interior state. At the boundaries, an attempted move outside the state space is
reflected, assuming reflected boundaries implies the chain stays put with probability \(1/2\) and moves inward with probability \(1/2\).
Thus, with states ordered as \(0,1,2,3,4\), the proposal matrix \(T=T(i,j)\) is
\[
T=
\begin{pmatrix}
\frac12 & \frac12 & 0 & 0 & 0 \\
\frac12 & 0 & \frac12 & 0 & 0 \\
0 & \frac12 & 0 & \frac12 & 0 \\
0 & 0 & \frac12 & 0 & \frac12 \\
0 & 0 & 0 & \frac12 & \frac12
\end{pmatrix}.
\]

\noindent Part B:

\noindent First compute the target probabilities:
\[
\pi(0)=\left(\frac34\right)^4=\frac{81}{256},\qquad
\pi(1)=4\left(\frac14\right)\left(\frac34\right)^3=\frac{27}{64},
\]
\[
\pi(2)=6\left(\frac14\right)^2\left(\frac34\right)^2=\frac{27}{128},\qquad
\pi(3)=4\left(\frac14\right)^3\left(\frac34\right)=\frac{3}{64},
\]
\[
\pi(4)=\left(\frac14\right)^4=\frac{1}{256}.
\]

\noindent Since the proposal is symmetric (\(T(i,j)=T(j,i)\)), the MH acceptance probability is
\[
\alpha(i,j)=\min\left(1,\frac{\pi(j)}{\pi(i)}\right).
\]
For neighboring moves this gives:
\[
\alpha(0,1)=1,\quad \alpha(1,0)=\frac34;\qquad
\alpha(1,2)=\frac12,\quad \alpha(2,1)=1;
\]
\[
\alpha(2,3)=\frac{2}{9},\quad \alpha(3,2)=1;\qquad
\alpha(3,4)=\frac{1}{12},\quad \alpha(4,3)=1.
\]

\noindent The MH transition probabilities for \(i\neq j\) are
\[
P(i,j)=T(i,j)\,\alpha(i,j),
\]
and the diagonal entries are
\[
P(i,i)=1-\sum_{j\neq i}P(i,j).
\]
Therefore,
\[
P(0,1)=\frac12\cdot 1=\frac12,\qquad P(0,0)=\frac12.
\]
\[
P(1,0)=\frac12\cdot\frac34=\frac38,\quad
P(1,2)=\frac12\cdot\frac12=\frac14,\quad
P(1,1)=1-\frac38-\frac14=\frac38.
\]
\[
P(2,1)=\frac12\cdot 1=\frac12,\quad
P(2,3)=\frac12\cdot\frac{2}{9}=\frac19,\quad
P(2,2)=1-\frac12-\frac19=\frac{7}{18}.
\]
\[
P(3,2)=\frac12\cdot 1=\frac12,\quad
P(3,4)=\frac12\cdot\frac{1}{12}=\frac{1}{24},\quad
P(3,3)=1-\frac12-\frac{1}{24}=\frac{11}{24}.
\]
\[
P(4,3)=\frac12\cdot 1=\frac12,\qquad P(4,4)=\frac12.
\]

\noindent Hence, with state order \(0,1,2,3,4\),
\[
P=
\begin{pmatrix}
\frac12 & \frac12 & 0 & 0 & 0 \\
\frac38 & \frac38 & \frac14 & 0 & 0 \\
0 & \frac12 & \frac{7}{18} & \frac19 & 0 \\
0 & 0 & \frac12 & \frac{11}{24} & \frac{1}{24} \\
0 & 0 & 0 & \frac12 & \frac12
\end{pmatrix}.
\]

\noindent Part C:

\noindent To show reversibility, we verify the detailed balance equations:
\[
\pi(i)P(i,j)=\pi(j)P(j,i)\qquad \text{for all } i,j.
\]
Since transitions only occur between neighbors, it suffices to check the pairs
\((0,1),(1,2),(2,3),(3,4)\).

\paragraph{Pair \((0,1)\).}
\[
\pi(0)P(0,1)=\frac{81}{256}\cdot\frac12=\frac{81}{512},\qquad
\pi(1)P(1,0)=\frac{27}{64}\cdot\frac38=\frac{81}{512}.
\]

\paragraph{Pair \((1,2)\).}
\[
\pi(1)P(1,2)=\frac{27}{64}\cdot\frac14=\frac{27}{256},\qquad
\pi(2)P(2,1)=\frac{27}{128}\cdot\frac12=\frac{27}{256}.
\]

\paragraph{Pair \((2,3)\).}
\[
\pi(2)P(2,3)=\frac{27}{128}\cdot\frac19=\frac{27}{1152},\qquad
\pi(3)P(3,2)=\frac{3}{64}\cdot\frac12=\frac{3}{128}=\frac{27}{1152}.
\]

\paragraph{Pair \((3,4)\).}
\[
\pi(3)P(3,4)=\frac{3}{64}\cdot\frac{1}{24}=\frac{3}{1536}=\frac{1}{512},\qquad
\pi(4)P(4,3)=\frac{1}{256}\cdot\frac12=\frac{1}{512}.
\]

\noindent Thus detailed balance holds for all possible moves, so the chain is reversible with respect to \(\pi\).
Finally, detailed balance implies stationarity: for each \(j\),
\[
\sum_i \pi(i)P(i,j)=\sum_i \pi(j)P(j,i)=\pi(j)\sum_i P(j,i)=\pi(j),
\]
so \(\pi P=\pi\). 

\subsection{Exercise 5.18}

\noindent A random variable $X$ has density function, defined up to proportionality,
\[
f(x) \propto e^{-(x-1)^2/2} + e^{-(x-4)^2/2}, \qquad 0 < x < 5.
\]
Implement a Metropolis--Hastings algorithm for simulating from the distribution.
Use your algorithm to approximate the mean and variance of $X$.



\noindent \textbf{Solution:}

\noindent \textbf{Input:}

\begin{lstlisting}[caption={Metropolis-Hastings algorithm}]
### Exercise 5.18

trials <- 1e6
sim <- numeric(trials)
sim[1] <- 2.45

# unnormalized target density on (0,5)
f_unnorm <- function(x) exp(-0.5 * (x - 1)^2) + exp(-0.5 * (x - 4)^2)

for (i in 2:trials) {
  old  <- sim[i - 1]
  prop <- runif(1, 0, 5)  # independence proposal q(prop) = Uniform(0,5)
  
  acc <- f_unnorm(prop) / f_unnorm(old)
  acc <- min(1, acc)      # acceptance probability
  
  sim[i] <- if (runif(1) < acc) prop else old
}

burn <- 10000
x <- sim[(burn + 1):trials]

mean(x)
var(x)
\end{lstlisting}

\noindent \textbf{Output:} Mean of 2.496 and Variance of 2.099.

\noindent \subsection{Exercise 5.19}

\noindent Implement the algorithm in Exercise 5.7 for $n = 50$ and $p = \tfrac{1}{4}$. 
Use your simulation to estimate 
\[
P(10 \le X \le 15),
\]
where $X$ has the given binomial distribution. 
Compare your estimate to the exact probability.


\noindent \textbf{Solution}

\noindent \textbf{Input:}

\begin{lstlisting}[caption={Metropolis-Hastings algorithm}]
### Exercise 5.19

trials <- 1e6
n <- 50
p <- 1/4

sim <- numeric(trials)

for (k in 1:trials) {
  state <- 0
  
  # run chain for 60 steps to be near stationarity
  for (i in 1:60) {
    y <- sample(0:n, 1)  # proposal: Uniform over {0,1,...,n}
    
    # log acceptance ratio:
    #   [C(n,y) p^y (1-p)^(n-y)] / [C(n,state) p^state (1-p)^(n-state)]
    log_acc <- (lchoose(n, y) - lchoose(n, state)) +
      (y - state) * log(p / (1 - p))
    
    acc <- min(1, exp(log_acc))
    
    if (runif(1) < acc) state <- y
  }
  
  sim[k] <- as.integer(state >= 10 && state <= 15)
}

mean(sim)  # estimate of P(10 <= X <= 15)

pbinom(15,n,p)-pbinom(9,n,p)
\end{lstlisting}

\noindent \textbf{Output:} Estimate = 0.671, Exact = .673

\end{document}